\section{Discusión}
\label{sec:discusion}

% ---------------------------------------------------------------
\subsection{El sesgo de densidad venía del llenado}
\label{sec:disc-sesgo}

El resultado central es que el sesgo de densidad que motivó varias calibraciones externas (condiciones C, D, E en la Tabla~\ref{tab:calibracion}) no era un defecto del mecanismo de la Memoria Asociativa Entrópica, sino una consecuencia de llenar mal las memorias.

Cuando la memoria se llenaba con el vector promedio de la clase repetido $N$ veces, todos los agentes acumulaban masas artificialmente iguales y los patrones de recall terminaban siempre en el mismo patrón (variedad $L_1 = 0.00$, Figura~\ref{fig:variedad-recall}).
Al reemplazar ese prototipo único por las instancias reales del conjunto de llenado, las masas de los agentes se igualaron por construcción: cada imagen aporta una contribución de igual peso y el gate de containment separa correctamente el 100\% de los casos de diagnóstico sin ninguna división adicional por la media de la memoria.

Dicho de forma simple, la teoría del modelo de Pineda y Morales~\cite{pineda2021entropic} estaba completa: bastaba alimentarla con los datos que el modelo requiere.
La lección práctica es que la forma en que se llena la memoria pesa más en el comportamiento del sistema que cualquier ajuste posterior de parámetros.

% ---------------------------------------------------------------
\subsection{La normalización necesaria es B1 en el directorio}
\label{sec:disc-b1}

A diferencia de las calibraciones que compensan síntomas, la lectura B1 del directorio transactivo no corrige un defecto: refleja cómo está construido el sistema.
Cuando los agentes especialistas acumulan números desiguales de interacciones (counts $[68, 54, 42]$), la puntuación cruda del directorio favorece de forma constante al agente más activo con independencia del contenido semántico de la consulta.
Dividir entre el conteo normaliza la evidencia a \textit{evidencia por interacción}, que es una medida justa para comparar agentes con historiales de exposición distintos.

Que B1 sea necesaria no implica que sea una corrección externa al modelo; al contrario, expresa que un agente puede tener más o menos experiencia en el sistema transactivo.
Las condiciones B2, C y D de la Tabla~\ref{tab:calibracion} intentan ajustar síntomas del mismo problema (la escala desigual de los scores) sin capturar la causa, lo que explica su menor desempeño (46.25\%, 33.75\% y 33.33\% respectivamente frente al 97.50\% de B1).

% ---------------------------------------------------------------
\subsection{Lectura desde la teoría de memoria transactiva de Wegner}
\label{sec:disc-wegner}

La teoría de memoria transactiva de Wegner~\cite{wegner1987transactive} dice que los grupos desarrollan un directorio compartido de quién sabe qué a través de la interacción repetida.
Los resultados de la Sección~\ref{sec:formacion} muestran este proceso con mucho detalle:

\begin{itemize}
    \item \textbf{Interacción mixta y rápida}: con presentación intercalada, el directorio cruza el umbral del 90\% en solo 14 interacciones.
    El grupo \textit{se conoce} muy pronto cuando la exposición es variada.
    \item \textbf{Segregación de experiencias}: con presentación bloqueada por dominio, la transición no ocurre hasta $k = 57$.
    Las experiencias homogéneas dentro de cada bloque aportan mucha información sobre un solo agente y casi ninguna sobre los demás, retrasando la formación del directorio colectivo.
    Esta dinámica recuerda los silos de conocimiento en organizaciones con poca rotación de roles.
    \item \textbf{Mediador sesgado}: el experimento con scoring crudo produce un grupo funcional (93.8\% de precisión final) pero con un directorio que distorsiona la especialización real: los counts $[73, 57, 34]$ sobreestiman la contribución del agente más activo frente a los counts equilibrados $[68, 54, 42]$ del sistema corregido.
    La lectura B1 compensa este sesgo en el agregado, pero la estructura interna del directorio queda distorsionada.
\end{itemize}

Estos resultados apuntan a que la calidad del directorio transactivo depende tanto del volumen de interacciones como de la distribución de esas interacciones entre los miembros del grupo.
Un sistema que solo maximiza la cantidad de interacciones sin mezclar bien las experiencias puede producir un directorio internamente sesgado que funciona bien en promedio pero falla al redistribuir tareas de forma equitativa.

% ---------------------------------------------------------------
\subsection{Limitaciones}
\label{sec:limitaciones}

\paragraph{(a) Cuantización por signo en la extracción semántica.}
La binarización por signo de los vectores fastText de 300 dimensiones descarta toda información de magnitud.
El margen discriminativo intra/inter-clase en el espacio binario es de solo 3.3 puntos porcentuales (60.4\% de acuerdo de signos intra-clase vs.\ 57.1\% inter-clase), lo que indica que las representaciones son similares entre clases en esta proyección.
Un efecto directo de esta limitación es la polisemia colapsada: los vectores de ``apple'' como fruta y ``apple'' como empresa tecnológica comparten el mismo vector fastText, lo que introduce ruido en el vocabulario de ConceptNet (etiquetas relacionadas con Apple Inc.\ aparecen en el vocabulario de la clase manzana-fruta).
La Figura~\ref{fig:espacio-semantico} muestra los solapamientos entre dominios en el espacio semántico 2D.

\begin{figure}[H]
    \centering
    \includegraphics[width=\textwidth]{espacio_semantico.png}
    \caption{Proyección bidimensional del espacio fastText de los labels de las tres clases.
    Los tres paneles muestran los solapamientos cross-domain que dificultan la discriminación por signo binario.}
    \label{fig:espacio-semantico}
\end{figure}

Las mitigaciones diseñadas incluyen: (i) usar las magnitudes de los componentes fastText como pesos por característica en lugar de binarizar, y (ii) una cuantización más fina en 16 niveles por componente para preservar gradientes de similitud.
Ambas alternativas son compatibles con la arquitectura EAM y son caminos directos para extender el sistema.

\paragraph{(b) Cobertura limitada del hemisferio visual con $\xi = 0$.}
Con el parámetro de generalización $\xi = 0$, el directorio visual solo acepta imágenes suficientemente similares a las ya registradas.
Esto produce una precisión perfecta sobre las imágenes aceptadas (100\%), pero una cobertura total de solo el 32.9\% en el conjunto de test.
El siguiente paso es probar varios valores de $\xi > 0$: valores pequeños de $\xi$ podrían aumentar la cobertura con un costo controlado en especificidad, dado que la tasa de falsos fue del 0\% en todo el rango de $N$ probado (Tabla~\ref{tab:capacidad}).

\paragraph{(c) Calidad de la reconstrucción visual.}
El decoder del autoencoder, entrenado con error cuadrático medio (MSE), produce reconstrucciones visualmente suaves que carecen de detalle de alta frecuencia.
Este es un efecto de la función de pérdida del autoencoder, no del sistema de memorias asociativas, y puede abordarse de forma independiente mediante funciones de pérdida perceptuales o arquitecturas generativas más expresivas.

\paragraph{(d) Aprendizaje limitado después de formar el directorio.}
En esta versión, una vez que se llenan las memorias de dominio y se forma el directorio,
los agentes dejan de registrar nuevo contenido semántico o visual.
La fase madura se usa para enrutar y recuperar con lo ya aprendido, no para seguir
almacenando labels, imágenes o asociaciones nuevas.
Por esa razón, el sistema todavía no cubre procesos como \emph{transactive encoding} y
\emph{transactive retrieval} en el sentido más amplio: no modela cómo el grupo incorpora
conocimiento nuevo ni cómo reorganiza sus memorias de contenido durante la operación
madura.
El directorio sí se evalúa como mecanismo de coordinación, pero el aprendizaje continuo
de los agentes queda fuera del alcance de este reporte.
